% Introduction: ~1 page.

An LLM agent that runs for more than a handful of turns has to remember
things. Tool outputs pile up, retrieved documents are consulted and
re-consulted, plans are revised, and sub-agents report back. The
dominant solution is to concatenate this history into the model's
context window and trust the attention mechanism to find what matters.
Even setting aside cost, this strategy fails for a measurable reason:
language models attend poorly to information in the middle of long
inputs \citep{liu2024lostmiddle}, and benchmarks such as LongBench
\citep{bai2024longbench} confirm that simply enlarging the window does
not close the gap.

A second line of work treats context as \emph{external memory}.
MemGPT \citep{packer2023memgpt} borrows the virtual-memory metaphor
from operating systems and pages content between the window and a
backing store. A-MEM \citep{xu2025amem} learns Zettelkasten-style
links between notes. Generative Agents \citep{park2023generative}
attach a memory stream with periodic reflection, and production
systems such as Mem0 \citep{chhikara2025mem0} assemble similar ideas
into deployable services. What these approaches share is an
\emph{opaque} interface: the agent emits a natural-language query, a
memory controller decides what to return, and the agent sees only the
result. The address space is hidden.

We argue that this opacity is the problem. An agent that is itself a
reasoning system should be able to reason about where its state
\emph{lives}, not only about what the state says. Concretely, an agent
should be able to ask: \emph{``Is there a tool that does X?''} and
receive a list of paths under \texttt{/tools}; it should be able to
distinguish \emph{``I need a summary of this document to decide if I
should read it''} from \emph{``I've decided; give me the full text.''}
Flat buffers, opaque retrievers, and learned links all collapse these
distinctions.

\paragraph{Our approach.} We propose \textsc{Structure}, a context
layer built on two primitives. First, every durable piece of state is
addressed by a hierarchical path and can be discovered by prefix and
glob queries --- the same operations a programmer uses on a
filesystem. Second, every entry is \emph{disclosable} at three
levels: a \textsc{glance} (one-line summary, $\leq 50$\,B), an
\textsc{overview} (structured metadata and tags, $\approx 1$\,KB),
and a \textsc{detail} (full content, optionally S3-backed). The
disclosure method lives on the entry itself, so an agent can scan
hundreds of glances in a single prompt and pay for full content only
where it is warranted. Mutations are emitted as events onto an
append-only log, giving us both real-time streaming to clients and
exact replay for evaluation.

\paragraph{Contributions.} Our main contributions are:
\begin{itemize}
    \item We formalise \emph{path-addressable multi-level disclosure}
    as a context-management primitive for LLM agents and contrast it
    with paging, retrieval, and learned-link memories
    (\S\ref{sec:related_work}).
    \item We introduce two mechanisms layered on top of the address
    space: (a) \emph{source evaluation}, where the agent rates each
    consulted entry and ratings are event-sourced into an aggregate
    that reorders future retrieval
    (\S\ref{sec:component_rating}); and (b) \emph{event-level
    garbage collection}, a per-type TTL policy that prunes transient
    events from the agent's working context on every step while
    preserving pinned turns and the append-only audit log
    (\S\ref{sec:component_gc}).
    \item We implement \textsc{Structure} as an event-sourced
    orchestration platform with a path-addressable context layer, three
    visibility scopes (user, workspace, global), and eleven first-class
    context tools; the system sustains sub-10\,ms retrieval over
    hundreds of entries via in-memory caches and HNSW indices
    (\S\ref{sec:method}).
    \item We evaluate the system on long-horizon agent tasks against
    flat-window, RAG, and MemGPT-style baselines, reporting accuracy,
    token cost, and end-to-end latency, and ablate the two new
    mechanisms to isolate their effect
    (\S\ref{sec:experiments}, \S\ref{sec:results}).
    \item We release the implementation, benchmark harness, and
    configurations so that the architecture can be reused and
    compared against future context-management designs.
\end{itemize}
